\honest{No, Really, I've Deceived Myself}{Eliezer Yudkowsky}{2009}

I recently talked for a few hours with a nominally Orthodox Jewish woman. She was highly intelligent and knew some of the evidence against her religion. For example, she knew that Mordecai, Esther, Haman and Vashti are not in the Persian records, and that there was an old legend about the gods Marduk, Ishtar, Humman and Vashti. She knows this, and she still celebrates Purim. \nb{The Persian records are not silent: the name Marduka, for officials of the Persian court, appears in over thirty Persepolis texts. The link to Marduk and Ishtar is a scholarly conjecture from the 1890s.}

``Most people like this'' will pretend to be too wise to talk to atheists, I say. I give no evidence.

What I learned is that you do not have to deceive yourself, so long as you believe you have. I call it ``belief in self-deception.''

In high school she thought she was an atheist, decided to act as if she believed in God, and, she told me, came to believe. \nb{Acting as if one believed until one does is the method Pascal recommended.}

``So far as I can tell,'' she is wrong about that. She always said ``I believe in God,'' never ``There is a God.'' She talked only about the benefits of believing. And she agreed that someone who just wanted the truth would not even invent God as a hypothesis. \nb{``I believe in God'' is the ordinary form of a profession of faith; in Jewish prayer, the thirteen lines of \textsc{Ani Ma'amin} each begin ``I believe with full faith.'' Her agreement about the hypothesis is the stronger evidence. All of it comes from one conversation, as the author recalls it.}

I think she does believe she has deceived herself. So she expects the benefits of believing in God, and that, ``I suppose,'' should produce ``much the same placebo effect'' as believing. I give no evidence for this. \nb{In ``Doublethink'' (2007) Yudkowsky had written that you ``might even believe you were happy and self-deceived; but you would not in fact be happy and self-deceived.'' The two claims can be reconciled, but the post does not do it.} This may explain why she defended ``I believe'' and was not interested in whether God exists.
